\section{Structural consequences and homological stability}
\label{sec:consequences}

Since \(S\) is finite, \(X(W,S)\) has finitely many cells, one for each
spherical subset of \(S\).  Theorem~\ref{thm:main} therefore gives a finite
classifying space for \(A\); its cellular chains and cochains compute the
integral group homology and cohomology.  The theorem also supplies the
\(K(\pi,1)\) hypothesis in the following results about torsion, centers,
positive monoids, and Boyd's homological-stability families.

\subsection{Torsion, centers, and monoid classifying spaces}

The Coxeter diagram has vertex set \(S\) and an edge between distinct
\(s,t\) exactly when \(m_{st}\ne2\).  Its connected components are the
irreducible Coxeter components, and a component is spherical when its
Coxeter group is finite.  Let \(A^+\) denote the positive Artin monoid:
it has the same standard generators and positive braid relations as
\(A\), interpreted as a monoid presentation.  For a monoid \(M\), write
\(BM\) for the geometric realization of the nerve of the one-object
category with endomorphism monoid \(M\).  A monoid homomorphism induces
a map of these classifying spaces.

\begin{corollary}[Torsion, center, and monoid consequences]
\label{cor:artin-structural}
Let \(A=A(W,S)\) be an Artin group with \(S\) finite, let \(A^+\) be its
positive Artin monoid, and let \(k\) be the number of irreducible
spherical components of the Coxeter diagram of \((W,S)\).  Then:
\begin{enumerate}[label=\textup{(\roman*)}]
  \item \(A\) is torsion-free.
  \item The center satisfies \(Z(A)\cong\Zbb^k\).
  \item The natural homomorphism \(\eta\colon A^+\to A\) induces a
    homotopy equivalence
    \[
      B\eta\colon BA^+\longrightarrow BA.
    \]
\end{enumerate}
\end{corollary}

\begin{proof}
When \(S=\varnothing\), both \(A\) and \(A^+\) are trivial and \(k=0\),
so all three assertions hold.  Assume that \(S\ne\varnothing\).

The complex \(X(W,S)\) has dimension at most \(|S|\), and it is a
\(K(A,1)\) by Theorem~\ref{thm:main} and
\eqref{eq:topology-fundamental-group}.  If \(A\) had a nonidentity
element of finite order, it would contain a subgroup \(P\cong C_p\)
for some prime \(p\).  The corresponding covering space
\(\widetilde X(W,S)/P\) would be a \(K(P,1)\) of dimension at most
\(|S|\).  This is impossible: the standard two-periodic resolution of
\(C_p\), whose cochain maps with trivial \(\Zbb\)-coefficients alternate
between zero and multiplication by \(p\), gives
\[
  H^{2q}(C_p;\Zbb)\cong\Zbb/p\Zbb\qquad(q\ge1),
\]
whereas the cohomology of a CW complex of dimension at most \(|S|\)
vanishes in higher degrees.  This proves (i).

Let \(S_1,\ldots,S_r\) be the connected components of the Coxeter
diagram.  All cross labels are \(2\), so the presentation gives
\[
  A\cong\prod_{j=1}^r A_{S_j},
\]
as in \eqref{eq:topology-parabolic-product}.  If \(S_j\) is spherical,
Deligne's center theorem for a finite Coxeter system with nonempty
connected diagram gives \(Z(A_{S_j})\cong\Zbb\)
\cite[Theorem~4.21]{Deligne1972}.  If \(S_j\) is not spherical, its
connectedness means that \(A_{S_j}\) has no spherical factor.
Theorem~\ref{thm:main} applied to the submatrix on \(S_j\) verifies the
\(K(\pi,1)\) hypothesis of Jankiewicz--Schreve's theorem
\cite[Theorem~3]{JankiewiczSchreve2022}, which gives
\(Z(A_{S_j})=\{1\}\).  Taking centers of the direct product proves (ii).

For (iii), Dobrinskaya's equivalence
\cite[Theorem~2 and the following paragraph]{Dobrinskaya2002}, reproved by Ozornova
\cite[Section~5, especially Conjecture~5.9 and Theorem~5.10]{Ozornova2013},
states that the natural map \(BA^+\to BA\) is a homotopy equivalence
exactly when the \(K(\pi,1)\) conjecture holds for \(A\).
Theorem~\ref{thm:main} supplies that condition.
\end{proof}

\subsection{Homological stability after adjoining a braid tail}

Boyd studies a fixed initial diagram with a growing type-\(A\) path.
We give the full matrix construction so that the indexing in her
stability range is explicit.

\begin{corollary}[Boyd stability for fixed-core braid-tail families]
\label{cor:boyd-stability}
Let \(A_1\) be the Artin group of a Coxeter matrix \(m_1\) on a finite
set \(\Sigma_1\), and choose a vertex \(\sigma_1\in\Sigma_1\).  For each
\(n\ge2\), adjoin new vertices \(\sigma_2,\ldots,\sigma_n\) and let
\(A_n\) be the Artin group of the Coxeter matrix \(m_n\) on
\(\Sigma_n=\Sigma_1\sqcup\{\sigma_2,\ldots,\sigma_n\}\).  Retain \(m_1\)
on \(\Sigma_1\), put diagonal entries equal to \(1\), and set the
remaining off-diagonal entries, symmetrically, by
\[
\begin{aligned}
 m_n(\sigma_j,\sigma_{j+1})&=3
     &&(1\le j<n),\\
 m_n(\sigma_j,\sigma_\ell)&=2
     &&(1\le j<\ell\le n,\ \ell-j\ge2),\\
 m_n(u,\sigma_j)&=2
     &&(u\in\Sigma_1\setminus\{\sigma_1\},\ 2\le j\le n).
\end{aligned}
\]
Let \(s_n\colon A_{n-1}\hookrightarrow A_n\) be the natural inclusion
for \(n\ge2\).  For every \(n\ge2\), every abelian group \(C\), used as
a constant coefficient group, and every integer \(i\ge0\), the induced map
\[
 (Bs_n)_*\colon H_i(BA_{n-1};C)\longrightarrow H_i(BA_n;C)
\]
is an isomorphism when \(i<n/2\) and a surjection when \(i=n/2\).
\end{corollary}

\begin{proof}
This is precisely Boyd's sequence of diagrams
\cite[Introduction, Definition~6.1 and Remark~6.2]{Boyd2020}.
Deleting \(\sigma_1\) from the initial diagram gives the diagram for
her \(A_0\), while \(\sigma_1,\ldots,\sigma_n\) form her type-\(A\)
subdiagram.  In particular, the new vertices have no other edges to
the fixed core, and our parameter \(n\) is Boyd's parameter.
The inclusions \(s_n\) are the standard parabolic inclusions.

Every \(\Sigma_n\) is finite, so Theorem~\ref{thm:main} proves the
\(K(\pi,1)\) conjecture for every \(A_n\) in this sequence.  This is
exactly the hypothesis of Boyd's Corollary~B, whose conclusion is the
displayed stabilization map and the stated constant-coefficient range.
\end{proof}

Here \(n\) counts the vertices in the growing type-\(A\) subdiagram,
not the total rank \(|\Sigma_1|+n-1\); the equality case \(i=n/2\) can
occur only when \(n\) is even.  The stability conclusion recorded here
is restricted to the displayed fixed-core construction and constant
coefficient systems.  It makes no assertion for arbitrary sequences
of Coxeter diagrams or for twisted coefficient systems.
